% GENERATED FILE. Edit canonical lessons, then run npm run generate:books.
% canonical-source-hash: fnv1a64:322ac7b1f48a0194
% canonical-lessons: ZH-C327-mingren, ZH-C327-yingxiong, ZH-C327-huangdi, ZH-C327-dazhang, ZH-C327-zhanzheng

\chapter{Heroes and Emperors}
\label{ch:zh-heroes-and-emperors}
\hlchaptermodality{327}

\begin{chapteropening}
\textbf{By the end of this chapter:} \emph{I can talk about famous people, heroes and emperors, and say two countries went to war.}

Complete the last lesson of chapter 327: I can talk about famous people, heroes and emperors, and say two countries went to war.
\end{chapteropening}

\section[\texorpdfstring{\zh{名人} (míngrén)}{名人 (míngrén)}]{\zh{名人} (míngrén) — famous person; well-known figure}

\label{lesson:ZH-C327-mingren}



\begin{quote}
Before the new one: say the Mandarin for to found something new, then the Mandarin for an ancient civilization.
\end{quote}

\subsection*{You'll want to know: \zh{名人}}
\textbf{\zh{名人}} — \emph{míngrén} — \textquotedblleft{}famous person; well-known figure\textquotedblright{}.

\begin{quote}
\textbf{\zh{他是一个世界名人}} \emph{tā shì yí ge shìjiè míngrén} — He's famous all over the world.
\end{quote}

\begin{cousinweb}
\textbf{\zh{名人}} (\emph{míngrén}) is built from \zh{名} 'a name' and \zh{人} 'a person', both inherited from Middle Chinese, \emph{mjieng} and \emph{nyin} (Guangyun, in Baxter's transcription), which glosses the first 'a name'.
\end{cousinweb}

\begin{grammarlens}[title={What you've built}]
One more word for telling what happened.
\end{grammarlens}

\subsection*{Your turn}
\begin{itemize}
\raggedright
  \item \emph{Say it:} \emph{míngrén}
  \item \emph{Say it:} \emph{míngrén}, once more
  \item \emph{Say it:} \emph{tā shì yí ge shìjiè míngrén} — He's famous all over the world.
  \item \emph{Recall:} say the Mandarin for to found something new, then the Mandarin for an ancient civilization, then say \emph{míngrén} again
\end{itemize}

\begin{tcolorbox}[breakable,colback=teal!4,colframe=teal!35!black,title={Before you move on}]
What does \zh{名人} mean? (\textquotedblleft{}Famous person; well-known figure\textquotedblright{}.) Say it once more.
\end{tcolorbox}
\section[yīngxióng]{yīngxióng — hero}

\label{lesson:ZH-C327-yingxiong}



\begin{quote}
Before the new one: say the Mandarin for an ancient civilization, then the Mandarin for a famous person.
\end{quote}

\subsection*{You'll want to know: yīngxióng}
\emph{yīngxióng} — \textquotedblleft{}hero\textquotedblright{}.

\begin{quote}
\emph{wǒ bàba shì wǒ de yīngxióng} — My dad is my hero.
\end{quote}

Its characters come later; for now it is a word for the ear.

\begin{cousinweb}
\emph{yīngxióng} is built from \emph{yīng} and \emph{xióng}, both inherited from Middle Chinese, \emph{'jaeng} and \emph{hjuwng} (Guangyun, in Baxter's transcription), which glosses the first 'a flower' and 'outstanding', and the second 'male' as against female.
\end{cousinweb}

\begin{grammarlens}[title={What you've built}]
One more word for telling what happened.
\end{grammarlens}

\subsection*{Your turn}
\begin{itemize}
\raggedright
  \item \emph{Say it:} \emph{yīngxióng}
  \item \emph{Say it:} \emph{yīngxióng}, once more
  \item \emph{Say it:} \emph{wǒ bàba shì wǒ de yīngxióng} — My dad is my hero.
  \item \emph{Recall:} say the Mandarin for an ancient civilization, then the Mandarin for a famous person, then say \emph{yīngxióng} again
\end{itemize}

\begin{tcolorbox}[breakable,colback=teal!4,colframe=teal!35!black,title={Before you move on}]
What does yīngxióng mean? (\textquotedblleft{}Hero\textquotedblright{}.) Say it once more.
\end{tcolorbox}
\section[huángdì]{huángdì — emperor}

\label{lesson:ZH-C327-huangdi}



\begin{quote}
Before the new one: say the Mandarin for a famous person, then the Mandarin for a hero.
\end{quote}

\subsection*{You'll want to know: huángdì}
\emph{huángdì} — \textquotedblleft{}emperor\textquotedblright{}.

\begin{quote}
\emph{tā shí suì de shíhou dāng le huángdì} — He became emperor when he was ten.
\end{quote}

Its characters come later; for now it is a word for the ear.

\begin{cousinweb}
\emph{huángdì} is built from \emph{huáng} and \emph{dì}, both inherited from Middle Chinese, \emph{hwang} and \emph{tejH} (Guangyun, in Baxter's transcription), which glosses the first 'a sovereign' and 'great', and calls the second the title of one who rules all under heaven.
\end{cousinweb}

\begin{grammarlens}[title={What you've built}]
One more word for telling what happened.
\end{grammarlens}

\subsection*{Your turn}
\begin{itemize}
\raggedright
  \item \emph{Say it:} \emph{huángdì}
  \item \emph{Say it:} \emph{huángdì}, once more
  \item \emph{Say it:} \emph{tā shí suì de shíhou dāng le huángdì} — He became emperor when he was ten.
  \item \emph{Recall:} say the Mandarin for a famous person, then the Mandarin for a hero, then say \emph{huángdì} again
\end{itemize}

\begin{tcolorbox}[breakable,colback=teal!4,colframe=teal!35!black,title={Before you move on}]
What does huángdì mean? (\textquotedblleft{}Emperor\textquotedblright{}.) Say it once more.
\end{tcolorbox}
\section[dǎzhàng]{dǎzhàng — to fight a war; to go into battle}

\label{lesson:ZH-C327-dazhang}



\begin{quote}
Before the new one: say the Mandarin for a hero, then the Mandarin for an emperor.
\end{quote}

\subsection*{You'll want to know: dǎzhàng}
\emph{dǎzhàng} — \textquotedblleft{}to fight a war; to go into battle\textquotedblright{}.

\begin{quote}
\emph{gǔdài de shíhou zhè liǎng ge guójiā chángcháng dǎzhàng} — In ancient times these two countries were often at war.
\end{quote}

Its characters come later; for now it is a word for the ear.

\begin{cousinweb}
\emph{dǎzhàng} is built from \zh{打} 'to strike' and \emph{zhàng} 'weapons', both inherited from Middle Chinese, \emph{taengX} and \emph{drjangH} (Guangyun, in Baxter's transcription), which glosses the first 'to hit' and the second 'arms'. Middle Chinese \emph{taengX} would not regularly give dǎ.
\end{cousinweb}

\begin{grammarlens}[title={What you've built}]
One more word for telling what happened.
\end{grammarlens}

\subsection*{Your turn}
\begin{itemize}
\raggedright
  \item \emph{Say it:} \emph{dǎzhàng}
  \item \emph{Say it:} \emph{dǎzhàng}, once more
  \item \emph{Say it:} \emph{gǔdài de shíhou zhè liǎng ge guójiā chángcháng dǎzhàng} — In ancient times these two countries were often at war.
  \item \emph{Recall:} say the Mandarin for a hero, then the Mandarin for an emperor, then say \emph{dǎzhàng} again
\end{itemize}

\begin{tcolorbox}[breakable,colback=teal!4,colframe=teal!35!black,title={Before you move on}]
What does dǎzhàng mean? (\textquotedblleft{}To fight a war; to go into battle\textquotedblright{}.) Say it once more.
\end{tcolorbox}
\section[zhànzhēng]{zhànzhēng — war}

\label{lesson:ZH-C327-zhanzheng}



\begin{quote}
Before the new one: say the Mandarin for an emperor, then the Mandarin for to go into battle.
\end{quote}

\subsection*{You'll want to know: zhànzhēng}
\emph{zhànzhēng} — \textquotedblleft{}war\textquotedblright{}.

\begin{quote}
\emph{zhànzhēng de shíhou yéye hái shì yí ge shàonián} — During the war my grandpa was still a teenager.
\end{quote}

Its characters come later; for now it is a word for the ear.

\begin{cousinweb}
\emph{zhànzhēng} is built from \emph{zhàn} and \emph{zhēng}, both inherited from Middle Chinese, \emph{tsyenH} and \emph{tsreang} (Guangyun, in Baxter's transcription), which glosses the first 'to fear' and the second 'to compete'. The sense 'to fight' for the first is not given there.
\end{cousinweb}

\begin{grammarlens}[title={What you've built}]
One more word for telling what happened.
\end{grammarlens}

\subsection*{Your turn}
\begin{itemize}
\raggedright
  \item \emph{Say it:} \emph{zhànzhēng}
  \item \emph{Say it:} \emph{zhànzhēng}, once more
  \item \emph{Say it:} \emph{zhànzhēng de shíhou yéye hái shì yí ge shàonián} — During the war my grandpa was still a teenager.
  \item \emph{Recall:} say the Mandarin for an emperor, then the Mandarin for to go into battle, then say \emph{zhànzhēng} again
\end{itemize}

\begin{tcolorbox}[breakable,colback=teal!4,colframe=teal!35!black,title={Before you move on}]
What does zhànzhēng mean? (\textquotedblleft{}War\textquotedblright{}.) Say it once more.
\end{tcolorbox}
